\section{Introduction}

Tool-augmented LLM agents expand the attack surface of language-model systems beyond what conventional input filtering can address. In a modern agent runtime, untrusted instructions may arrive not only through direct user prompts but also through retrieved web content, tool-returned text, intermediate prompt construction steps, outbound message generation, and modifications to local control files. Once an agent acquires the ability to browse the web, invoke shell commands, persist intermediate artifacts, and emit externally visible output, security risks---prompt override, tool abuse, credential leakage, control-file tampering---become distributed across the entire runtime rather than confined to a single input boundary \cite{securityofaiagents2024,asb2024,owaspllmtop102025,mitreatlas2023}.

Existing defenses, however, still operate predominantly at one or two checkpoints, typically by classifying incoming prompts or filtering final outputs. While these boundary-focused approaches remain useful, they are structurally mismatched to the execution model of tool-using agents. Indirect prompt injection, for instance, may arrive through fetched pages or tool-returned results long after the original user prompt has passed inspection \cite{greshake2023,llamafirewall2025,indirectpromptfirewalls2025}. Low-grade suspicious signals may accumulate across multiple conversational turns before crossing a meaningful risk threshold. Tool misuse and outbound credential exfiltration often demand policy enforcement rather than text classification alone \cite{safe-tool-use2026}. For agent gateways, the central challenge is therefore not merely whether a single text string appears malicious, but how security controls can be distributed coherently across a multi-stage runtime that interleaves prompts, tool invocations, mutable state, and operator workflows.

This challenge is as much an engineering problem as it is a modeling problem. In practice, runtime defenses must be deployable without invasive modifications to the host framework, maintainable as that framework evolves, and observable by the operators who depend on them. Security logic that requires forking the upstream codebase becomes a maintenance burden with every upstream release. Static policies that cannot be reloaded, inspected, or verified at runtime are difficult to operate under real conditions. Detection-only approaches are likewise insufficient for production gateways, where operators need explicit enforcement points, auditable security events, and recovery workflows for handling policy updates and exceptions. This emphasis on deployment-time governance is also consistent with broader AI risk-management guidance that treats safety as a socio-technical operational problem rather than a model-only property \cite{nistairmf2023}. A practical agent defense system must therefore function as an operational runtime layer---not as a benchmark-only detector.

We present OpenClaw PRISM, a zero-fork runtime security layer for OpenClaw-based tool-augmented LLM agent gateways. PRISM combines an in-process plugin with optional sidecar services and distributes enforcement across ten lifecycle hooks covering message ingress, prompt construction, tool execution, tool-result persistence, outbound messaging, sub-agent spawning, and gateway startup. Rather than introducing a novel detection model, PRISM integrates several complementary defense mechanisms within a unified runtime architecture:

\begin{itemize}[leftmargin=1.5em]
  \item \textbf{Lifecycle-wide interception} across ten hooks that span the full agent interaction cycle.
  \item \textbf{A hybrid scanning pipeline} that applies fast heuristic scoring first and can escalate suspicious tool results to an LLM-assisted classifier.
  \item \textbf{Session-level and conversation-level risk accumulation} with time-decaying state and graduated response thresholds.
  \item \textbf{Policy-enforced controls} over tool invocations, protected paths, private networks, domain tiers, and outbound secret patterns.
  \item \textbf{A tamper-evident audit and operations plane} with chained audit records, integrity verification, hot-reloadable policy configuration, and dashboard allow workflows.
\end{itemize}

The goal is deployable runtime defense: a system that can be attached to an existing gateway, supervised as long-running services, audited after the fact, and tuned by operators---all without forking the upstream framework.

This paper makes the following contributions:

\begin{enumerate}[leftmargin=1.5em]
  \item \textbf{Zero-fork runtime defense architecture.} We design and implement a runtime defense layer for OpenClaw-based agent gateways that combines an in-process plugin with optional sidecar services, avoiding any modification to upstream gateway code.
  \item \textbf{Lifecycle-wide enforcement with staged risk response.} We implement enforcement across ten runtime hooks together with a hybrid heuristic-first, LLM-assisted scanning path and thresholded risk accumulation that supports graduated responses rather than one-shot classification.
  \item \textbf{Unified tool and network governance.} We unify tool governance, protected-path enforcement, private-network and domain controls, and outbound secret filtering within a single runtime security layer for tool-using agents.
  \item \textbf{Tamper-evident audit and operational plane.} We provide an operational security plane that includes chained audit records with integrity verification, component health probing, hot-reloadable policy configuration, and dashboard allow workflows.
  \item \textbf{Evaluation methodology and benchmark artifacts.} We outline an evaluation methodology and accompanying benchmark artifact structure for measuring security effectiveness, false positives, layer contribution, runtime overhead, and operational recoverability in an agent-runtime setting.
\end{enumerate}

The remainder of this paper is organized as follows. Section~2 defines the agent runtime model and threat model. Section~3 presents the PRISM architecture and its core defense mechanisms. Section~4 describes the implementation and deployment model. Section~5 outlines the evaluation methodology and reports current preliminary benchmark results. Section~6 positions PRISM relative to related work, and Section~7 discusses limitations and future directions.
